% ============================================================================
% TECHNICAL APPENDIX 3A: Machine Learning Foundations for HITL Design
% The mathematics of training, generalization, and distribution shift
% ============================================================================
% Source: /markdown/ch3/Ch3A_final.md
% ============================================================================

\chapter*{Technical Appendix 3A: Machine Learning Foundations for HITL Design}
\addcontentsline{toc}{chapter}{Technical Appendix 3A: Machine Learning Foundations for HITL Design}

\begin{chapterquote}{Formal foundations for Chapter 3}
Chapter~3 introduced machine learning as a three-step process of gathering examples, learning patterns, and generalizing. This appendix provides the mathematical framework underlying those concepts: empirical risk minimization, the bias-variance tradeoff, generalization bounds, and distribution shift.
\end{chapterquote}

% ============================================================================
\section{Empirical Risk Minimization}
% ============================================================================

The core framework for supervised machine learning is \term{empirical risk minimization} (ERM). Given training data $\mathcal{D} = \{(x_i, y_i)\}_{i=1}^n$, we find parameters $\theta$ that minimize:

\[
\hat{R}(\theta) = \frac{1}{n} \sum_{i=1}^{n} \mathcal{L}\!\left(f_\theta(x_i),\, y_i\right)
\]

where $\mathcal{L}$ is a loss function (e.g., cross-entropy for classification, mean squared error for regression) and $f_\theta$ is the model.

The \textbf{true risk} we actually care about is:

\[
R(\theta) = \mathbb{E}_{(x,y) \sim P}\!\left[\mathcal{L}\!\left(f_\theta(x),\, y\right)\right]
\]

The \term{generalization gap} is:

\[
R(\theta) - \hat{R}(\theta) \geq 0
\]

This gap is always non-negative (empirical risk is optimistically biased) and tends to grow with model complexity relative to dataset size.

% ============================================================================
\section{The Bias-Variance Tradeoff}
% ============================================================================

For a regression model, the expected prediction error decomposes as:

\[
\mathbb{E}\!\left[(y - \hat{f}(x))^2\right]
= \underbrace{\mathrm{Bias}^2\!\left(\hat{f}(x)\right)}_{\text{Systematic error}}
+ \underbrace{\mathrm{Var}\!\left(\hat{f}(x)\right)}_{\text{Sensitivity to training data}}
+ \underbrace{\sigma^2}_{\text{Irreducible noise}}
\]

\begin{itemize}
    \item \textbf{Bias:} Error from the model not capturing the true relationship (underfitting)
    \item \textbf{Variance:} Error from sensitivity to fluctuations in training data (overfitting)
    \item \textbf{$\sigma^2$:} Aleatoric uncertainty~--- cannot be reduced
\end{itemize}

The bias-variance tradeoff: reducing bias (using a more complex model) typically increases variance (model fits noise), and vice versa.

\begin{examplebox}[HITL Design Implication of Bias vs.\ Variance]
\begin{itemize}
    \item \textbf{High-variance models} benefit most from human-in-the-loop review because their errors are unpredictable~--- concentrated at unusual inputs with no systematic pattern.
    \item \textbf{High-bias models} benefit most from systematic data collection for the classes of inputs where bias is highest.
\end{itemize}
\end{examplebox}

% ============================================================================
\section{Generalization Bounds}
% ============================================================================

\subsection{VC Dimension Bound}

For a hypothesis class $\mathcal{H}$ with VC dimension $d_{\mathrm{VC}}$, with probability $1 - \delta$:

\[
R(\theta) \leq \hat{R}(\theta) + \sqrt{\frac{8}{n}\!\left(d_{\mathrm{VC}} \cdot \ln\frac{2n}{d_{\mathrm{VC}}} + \ln\frac{4}{\delta}\right)}
\]

The generalization gap shrinks at rate $\mathcal{O}\!\left(\sqrt{d_{\mathrm{VC}}/n}\right)$. More complex models (higher $d_{\mathrm{VC}}$) need more data to close the gap.

\subsection{PAC-Bayes Bound}

A tighter bound for stochastic classifiers:

\[
\mathbb{E}_{\theta \sim Q}[R(\theta)] \leq \mathbb{E}_{\theta \sim Q}[\hat{R}(\theta)] + \sqrt{\frac{\mathrm{KL}(Q \| P) + \ln(n/\delta)}{2n}}
\]

where $P$ is the prior and $Q$ is the posterior distribution over parameters.

% ============================================================================
\section{Distribution Shift}
% ============================================================================

\subsection{Types of Shift}

\paragraph{Covariate shift.} $P(x)$ changes but $P(y \mid x)$ stays the same:
\[
P_{\mathrm{train}}(x) \neq P_{\mathrm{test}}(x), \quad \text{but} \quad P_{\mathrm{train}}(y \mid x) = P_{\mathrm{test}}(y \mid x)
\]
\textit{Example:} a spam filter trained on emails from 2020, deployed in 2025~--- email styles changed but spam is still spam.

\paragraph{Label shift.} $P(y)$ changes but $P(x \mid y)$ stays the same:
\[
P_{\mathrm{train}}(y) \neq P_{\mathrm{test}}(y), \quad \text{but} \quad P_{\mathrm{train}}(x \mid y) = P_{\mathrm{test}}(x \mid y)
\]
\textit{Example:} fraud rate changes seasonally; the appearance of fraud is similar but its base rate changes.

\paragraph{Concept drift.} The relationship $P(y \mid x)$ changes:
\[
P_{\mathrm{train}}(y \mid x) \neq P_{\mathrm{test}}(y \mid x)
\]
\textit{Example:} a hiring model trained before a policy change; the definition of ``qualified'' shifts.

\paragraph{Wildcard.} A completely new feature space:
\[
\exists\, x \in \mathcal{X}_{\mathrm{test}} : p_{\mathrm{train}}(x) \approx 0
\]
\textit{Example:} a vision model encountering medical images after being trained on natural photographs.

\subsection{Detecting Distribution Shift}

\paragraph{Maximum Mean Discrepancy (MMD).}

\[
\mathrm{MMD}^2(P, Q) = \left\|\mathbb{E}_{x \sim P}[\phi(x)] - \mathbb{E}_{x \sim Q}[\phi(x)]\right\|^2_{\mathcal{H}}
\]

For an RBF kernel $k(x, x') = \exp\!\left(-\|x - x'\|^2 / 2\sigma^2\right)$:

\begin{lstlisting}[style=python, caption={MMD estimation between training and test distributions}]
def compute_mmd(X_train, X_test, sigma=1.0):
    """Estimate MMD between training and test distributions."""
    def rbf_kernel(X, Y):
        dists = cdist(X, Y, 'sqeuclidean')
        return np.exp(-dists / (2 * sigma**2))

    K_tt = rbf_kernel(X_train, X_train)
    K_ee = rbf_kernel(X_test,  X_test)
    K_te = rbf_kernel(X_train, X_test)

    return K_tt.mean() - 2*K_te.mean() + K_ee.mean()
\end{lstlisting}

\paragraph{Simple monitoring approach for production HITL.}

\begin{lstlisting}[style=python, caption={Distribution monitor with KS-test drift detection}]
class DistributionMonitor:
    def __init__(self, reference_predictions, window_size=1000):
        self.reference   = np.array(reference_predictions)
        self.window_size = window_size
        self.recent      = deque(maxlen=window_size)

    def update(self, new_prediction):
        self.recent.append(new_prediction)

    def drift_score(self):
        if len(self.recent) < self.window_size:
            return None
        recent = np.array(self.recent)
        # KS test p-value (low p = significant drift)
        stat, p_value = ks_2samp(self.reference, recent)
        return stat, p_value
\end{lstlisting}

% ============================================================================
\section{The PAC Learning Framework}
% ============================================================================

Probably Approximately Correct (PAC) learning formalizes when learning is feasible.

A concept class $\mathcal{C}$ is PAC learnable if there exists an algorithm $\mathcal{A}$ such that for any $\varepsilon, \delta > 0$ and any distribution $P$ over $\mathcal{X} \times \mathcal{Y}$, given $m \geq \mathrm{poly}(1/\varepsilon,\, 1/\delta,\, n,\, \mathrm{size}(c))$ examples, with probability $\geq 1 - \delta$:

\[
R\!\left(\mathcal{A}(\mathcal{D}_m)\right) \leq \varepsilon
\]

\begin{keyconcept}[PAC Theory and Rare-Category HITL]
Learning requires a finite number of labeled examples proportional to $1/\varepsilon$. For very rare categories (low prior probability), the sample required to learn them grows rapidly. Rare-category detection is a natural HITL application: human review of uncertain cases on rare categories is often the only practical way to achieve acceptable error rates.
\end{keyconcept}

% ============================================================================
\section{Exercises}
% ============================================================================

\begin{Exercise}[Distribution Shift Diagnosis]
A spam filter achieves 99.2\% accuracy on its test set but 96.1\% accuracy in deployment. The test set was drawn from the same email corpus as training. Identify: (a) what type of distribution shift is most likely; (b) how you would diagnose it; (c) what HITL intervention would help most.
\end{Exercise}

\begin{Exercise}[Bias-Variance and Fine-Tuned Models]
Using the bias-variance decomposition, explain why a large pre-trained language model fine-tuned on a small domain-specific dataset might be riskier in HITL deployment than a simpler model trained on a larger domain-specific dataset. What monitoring would you add?
\end{Exercise}

\begin{Exercise}[MMD Monitoring]
Implement a simple version of maximum mean discrepancy monitoring in Python. Use it to detect when a batch of new predictions has drifted from a reference batch. What threshold for MMD would you use to trigger a human review audit?
\end{Exercise}

\begin{Exercise}[Irreducible Error and HITL]
The irreducible error $\sigma^2$ in the bias-variance decomposition is analogous to aleatoric uncertainty. Explain in technical terms why human-in-the-loop intervention cannot reduce this component of error, and what a HITL system should do when it detects a case with high aleatoric uncertainty.
\end{Exercise}

% ============================================================================
\section{Recommended Reading}
% ============================================================================

\subsection*{Foundational}
\begin{itemize}
    \item Vapnik, V. (1998). \textit{Statistical Learning Theory}. Wiley.
    \item Shalev-Shwartz, S., \& Ben-David, S. (2014). \textit{Understanding Machine Learning: From Theory to Algorithms}. Cambridge University Press.
\end{itemize}

\subsection*{Distribution Shift}
\begin{itemize}
    \item Quinonero-Candela, J., et al.\ (Eds.). (2009). \textit{Dataset Shift in Machine Learning}. MIT Press.
    \item Sugiyama, M., \& Kawanabe, M. (2012). \textit{Machine Learning in Non-Stationary Environments}. MIT Press.
    \item Gretton, A., et al.\ (2012). A kernel two-sample test. \textit{Journal of Machine Learning Research}, 13, 723--773.
\end{itemize}

\subsection*{PAC Learning}
\begin{itemize}
    \item McAllester, D. (1999). PAC-Bayesian model averaging. \textit{COLT}.
\end{itemize}

\bigskip
\begin{center}
\rule{0.5\textwidth}{0.4pt}
\end{center}
\bigskip

\noindent\textit{This appendix provides the mathematical foundations for empirical risk minimization, the bias-variance tradeoff, generalization bounds, and distribution shift as introduced in Chapter~3. The worked exercises connect theory to practical HITL monitoring and intervention design.}
